This project will develop a geometry-aware, foundation-model–driven,
agentic AI framework for the rapid prediction, optimization, and control of
thermal–fluid transport in DOE-relevant energy systems. The proposed work
directly aligns with the chosen focus topic on Artificial Intelligence–enabled
design and control of complex energy systems by integrating high-fidelity
simulation data, advanced machine learning architectures, and closed-loop
optimization workflows.
%
Specifically, the project will: \\[-3ex]
\begin{enumerate}
\item \textbf{Develop a unified data representation and baseline surrogate}
that integrates heterogeneous experimental and simulation datasets from
\textbf{Nek5000/RS} and \textbf{MFEM} into a common geometry-aware latent
space suitable for large-scale foundation-model training.
\\[-3ex]

\item \textbf{Construct and fine-tune a geometry-aware foundation model}
that enables cross-configuration prediction of thermal–fluid quantities
(e.g., heat-transfer coefficients and pressure drop) with target error
$< \; 5\%$ across unseen geometries/parameters.
\\[-3ex]

\item \textbf{Enable latent-space optimization and closed-loop design}
through integration of Bayesian optimization, derivative-free methods,
and reinforcement learning, reducing the number of required high-fidelity
simulations by at least $5$--$10\times$ relative to brute-force CFD-based
design exploration.
\\[-3ex]

\item \textbf{Deploy an agentic AI orchestration framework} that performs
adaptive sampling and multi-objective optimization, demonstrating
end-to-end automated design workflows with $5$--$10\times$ acceleration
in design-cycle time for representative heat-transfer systems.
\\[-3ex]
\end{enumerate}
   These objectives will amplify in-house expertise in thermal-hydraulics
device design for fission and fusion systems.
   The AI advantage arises from the ability to replace large ensembles of
computationally expensive simulations with learned, physics-informed surrogate
models that generalize across geometries and operating conditions. 
  The amortized training will yield design-cycle cost, accelerate innovation
timelines, and enable exploration of previously inaccessible design spaces.

